\chapter{Hasil Kali Skalar pada Bidang}\label{ch:g11:scal}

\omterm{def:g11:scal:dot}{Hasil kali skalar} melekatkan sebuah \emph{bilangan} pada sepasang \omterm{def:g11:vect:vector}{vektor}, dan
bilangan itu melihat geometri: nilainya lenyap tepat untuk arah yang saling
tegak lurus, dan bilangan itu mengukur sudut serta panjang. Inilah mesin di
balik aturan \omterm{def:g11:trigo:cossin}{kosinus}, \omterm{def:g10:algebra:equation}{persamaan} lingkaran, dan perhitungan ketegaklurusan.
Perkakas yang sama, di ruang, menjadi pokok \cref{ch:g12:space}.

\section{Definisi dan perhitungan pertama}

\begin{definition}[Hasil kali skalar]\label{def:g11:scal:dot}
\emph{Hasil kali skalar}\index{hasil kali skalar} dua \omterm{def:g11:vect:vector}{vektor} $\vec u$
dan $\vec v$ adalah bilangan
\[
\vec u \cdot \vec v
= \tfrac12\left(\norm{\vec u + \vec v}^2 - \norm{\vec u}^2
- \norm{\vec v}^2\right),
\]
dengan $\norm{\vec u}$ menyatakan panjang (\emph{norma}\index{norma}) $\vec u$.
\end{definition}

\begin{theorem}[Empat wajah hasil kali skalar]\label{thm:g11:scal:faces}
Misalkan $\vec u\,(x, y)$ dan $\vec v\,(x', y')$ \omterm{def:g11:vect:vector}{vektor} taknol, serta
$\theta$ sudut di antara keduanya. Maka
\begin{enumerate}
  \item \emph{(\omterm{def:g10:coordgeom:system}{koordinat})} $\vec u \cdot \vec v = xx' + yy'$;
  \item \emph{(\omterm{def:g11:trigo:cossin}{kosinus})} $\vec u \cdot \vec v =
        \norm{\vec u}\,\norm{\vec v}\cos\theta$;
  \item \emph{(proyeksi)} $\vec u \cdot \vec v = \vec u \cdot \vec v'$,
        dengan $\vec v'$ proyeksi \omterm{def:g11:scal:orthogonal}{ortogonal} $\vec v$ pada
        arah $\vec u$; jika $\vec u = \vect{OA}$ dan
        $\vec v = \vect{OB}$, maka
        $\vec u \cdot \vec v = \overline{OA} \times \overline{OH}$, yaitu
        hasil kali panjang bertanda, dengan $H$ kaki
        garis tegak lurus dari $B$ ke garis $(OA)$;
  \item \emph{(\omterm{def:g11:scal:dot}{norma})} rumus pendefinisi pada \cref{def:g11:scal:dot}.
\end{enumerate}
\end{theorem}

\begin{proof}[Bukti butir 1 dan 2]
\emph{1.} Dalam \omterm{def:g10:coordgeom:system}{koordinat}, $\norm{\vec u}^2 = x^2 + y^2$ (Pythagoras) dan
$\vec u + \vec v$ berkoordinat $(x + x', y + y')$, sehingga
\begin{align*}
2\,\vec u \cdot \vec v
&= (x + x')^2 + (y + y')^2 - (x^2 + y^2) - (x'^2 + y'^2)\\
&= 2xx' + 2yy',
\end{align*}
setelah kuadratnya dijabarkan lalu saling dihapuskan.

\emph{2.} Pilihlah \omterm{def:g10:coordgeom:system}{sistem koordinatnya} sehingga $\vec u$ menunjuk sepanjang
sumbu $x$ positif: $\vec u\,(\norm{\vec u}, 0)$ dan
$\vec v\,(\norm{\vec v}\cos\theta,\ \norm{\vec v}\sin\theta)$
(\cref{ch:g11:trigo}). Rumus 1 lalu memberi
$\vec u \cdot \vec v = \norm{\vec u}\,\norm{\vec v}\cos\theta + 0$.
Rumus 3 terbaca dari perhitungan yang sama: proyeksi $\vec v$ pada
sumbu $x$ berpanjang bertanda $\norm{\vec v}\cos\theta =
\overline{OH}$.
\end{proof}

\begin{omfigure}
\begin{tikzpicture}[scale=1.25]
  \coordinate (O) at (0,0);
  \coordinate (A) at (3.2,0);
  \coordinate (B) at (55:2.4);
  \coordinate (H) at (1.3767,0);
  \draw[-Stealth, omDef, thick] (O) -- (A)
    node[below, font=\small, pos=0.85] {$\vec u$};
  \draw[-Stealth, omThm, thick] (O) -- (B)
    node[above left, font=\small, pos=0.8] {$\vec v$};
  \draw[black, dashed, thin] (B) -- (H);
  \draw[thin] (H) ++(-0.18,0) -- ++(0,0.18) -- ++(0.18,0);
  \draw[omProp, very thick] (O) -- (H);
  \fill (H) circle (1.2pt) node[below, font=\small] {$H$};
  \node[below left, font=\small] at (O) {$O$};
  \draw[->,omThm,thin] (0.55,0) arc (0:55:0.55);
  \node[omThm, font=\small] at (27:0.78) {$\theta$};
\end{tikzpicture}

{\small Tampilan proyeksinya: $\vec u \cdot \vec v = \overline{OA} \times
\overline{OH}$, dengan $H$ kaki garis tegak lurus dari ujung
$\vec v$. Ruas jingga $OH$ berpanjang bertanda
$\norm{\vec v}\cos\theta$.}
\end{omfigure}

\begin{example}\label{ex:g11:scal:compute}
$\vec u\,(3, 1)$ dan $\vec v\,(2, -4)$:
$\vec u \cdot \vec v = 6 - 4 = 2$. \omterm{def:g11:scal:dot}{Normanya}
$\norm{\vec u} = \sqrt{10}$ dan $\norm{\vec v} = \sqrt{20}$, sehingga
$\cos\theta = \frac{2}{\sqrt{10}\sqrt{20}} = \frac{2}{10\sqrt2}
= \frac{\sqrt2}{10}$: sudutnya sedikit kurang dari $82^\circ$.
\end{example}

\begin{proposition}[Aturan aljabar]\label{prop:g11:scal:rules}
Untuk semua \omterm{def:g11:vect:vector}{vektor} dan semua $\lambda \in \R$:
\[
\vec u \cdot \vec v = \vec v \cdot \vec u, \qquad
\vec u \cdot (\vec v + \vec w) = \vec u \cdot \vec v + \vec u \cdot
\vec w, \qquad
(\lambda \vec u)\cdot \vec v = \lambda\,(\vec u \cdot \vec v),
\]
\[
\vec u \cdot \vec u = \norm{\vec u}^2,
\qquad
\norm{\vec u + \vec v}^2 = \norm{\vec u}^2 + 2\,\vec u\cdot\vec v +
\norm{\vec v}^2 .
\]
\end{proposition}

\begin{proof}
Semuanya menyusul dari rumus \omterm{def:g10:coordgeom:system}{koordinatnya}: misalnya
$\vec u \cdot (\vec v + \vec w) = x(x' + x'') + y(y' + y'') =
(xx' + yy') + (xx'' + yy'')$. Identitas terakhirnya \omterm{def:g10:algebra:expand}{menjabarkan}
$\norm{\vec u + \vec v}^2 = (\vec u + \vec v)\cdot(\vec u + \vec v)$
dengan ketiga aturan pertamanya --- \omterm{def:g11:scal:dot}{hasil kali skalar} berperilaku persis
seperti hasil kali biasa, sehingga identitas aljabar yang lazim berlaku pula
untuk \omterm{def:g11:vect:vector}{vektor}.
\end{proof}

\section{Ketegaklurusan}

\begin{definition}[Vektor ortogonal]\label{def:g11:scal:orthogonal}
Dua \omterm{def:g11:vect:vector}{vektor} disebut \emph{ortogonal}\index{vektor ortogonal}, ditulis
$\vec u \perp \vec v$, bila $\vec u \cdot \vec v = 0$.
\end{definition}

\begin{remark}
Untuk \omterm{def:g11:vect:vector}{vektor} taknol, rumus 2 pada \cref{thm:g11:scal:faces} memperlihatkan
$\vec u \cdot \vec v = 0 \iff \cos\theta = 0 \iff$ arahnya saling
tegak lurus. \omterm{def:g11:vect:vector}{Vektor} nol \omterm{def:g11:scal:orthogonal}{ortogonal} terhadap apa saja.
\end{remark}

\begin{example}
$\vec u\,(3, 2)$ dan $\vec v\,(-4, 6)$: $\vec u \cdot \vec v = -12 + 12 =
0$: jadi \omterm{def:g11:scal:orthogonal}{ortogonal}. Secara umum, $\vec u\,(a, b)$ selalu \omterm{def:g11:scal:orthogonal}{ortogonal} terhadap
$\vec v\,(-b, a)$ --- yaitu rotasi seperempat putaran.
\end{example}

\section{Aturan kosinus}

\begin{theorem}[Aturan kosinus]\label{thm:g11:scal:cosines}
\index{aturan kosinus}
Pada segitiga $ABC$ dengan sisi $a = BC$, $b = CA$, $c = AB$ dan sudut
$\widehat A$ di titik sudut $A$:
\[
a^2 = b^2 + c^2 - 2bc\,\cos\widehat A .
\]
\end{theorem}

\begin{proof}
Tulislah $\vect{BC} = \vect{BA} + \vect{AC} = \vect{AC} - \vect{AB}$ lalu
jabarkan dengan \cref{prop:g11:scal:rules}:
\[
a^2 = \norm{\vect{BC}}^2
= \norm{\vect{AC}}^2 - 2\,\vect{AC}\cdot\vect{AB} + \norm{\vect{AB}}^2
= b^2 + c^2 - 2\,\vect{AB}\cdot\vect{AC}.
\]
Menurut rumus kosinusnya,
$\vect{AB}\cdot\vect{AC} = c\,b\,\cos\widehat A$.
\end{proof}

\begin{remark}
Ketika $\widehat A$ siku-siku, $\cos\widehat A = 0$ dan aturan
kosinusnya menyusut menjadi teorema Pythagoras $a^2 = b^2 + c^2$: jadi aturan
itu adalah Pythagoras dengan suku koreksi untuk sudut yang tidak siku-siku.
\end{remark}

\begin{omfigure}
\begin{tikzpicture}[scale=1.15]
  \coordinate (A) at (0,0);
  \coordinate (B) at (4,0);
  \coordinate (C) at (1.2,2.2);
  \draw[omDef, thick] (A) -- (B) -- (C) -- cycle;
  \node[below left, font=\small] at (A) {$A$};
  \node[below right, font=\small] at (B) {$B$};
  \node[above, font=\small] at (C) {$C$};
  \node[below, font=\small] at (2,0) {$c$};
  \node[above left, font=\small] at (0.6,1.1) {$b$};
  \node[above right, font=\small] at (2.6,1.1) {$a$};
  \draw[->,omThm,thin] (0.7,0) arc (0:61.4:0.7);
  \node[omThm, font=\small] at (30:1.0) {$\widehat A$};
\end{tikzpicture}

{\small Aturan \omterm{def:g11:trigo:cossin}{kosinus}: mengetahui dua sisi $b$, $c$ dan sudut di
antaranya sudah menentukan sisi ketiganya, $a$.}
\end{omfigure}

\begin{example}\label{ex:g11:scal:cosines}
Dua sisi sebuah segitiga berukuran $b = 5$ dan $c = 8$, dan mengapit sudut
$\widehat A = \frac\pi3$. Maka
\[
a^2 = 25 + 64 - 2 \times 5 \times 8 \times \frac12 = 89 - 40 = 49,
\]
sehingga $a = 7$.
\end{example}

\section{Penerapan: lingkaran dan garis tegak lurus}

\begin{proposition}[Lingkaran dengan diameter tertentu]\label{prop:g11:scal:diameter}
Sebuah titik $M$ terletak pada lingkaran berdiameter $[AB]$ jika dan hanya jika
\[
\vect{MA} \cdot \vect{MB} = 0 .
\]
\end{proposition}

\begin{proof}
Jika $M \neq A, B$, syarat itu mengatakan bahwa sudut di $M$ pada segitiga
$AMB$ siku-siku, dan titik yang memandang sebuah ruas garis dengan sudut
siku-siku tepat adalah titik pada lingkaran yang berdiameter ruas itu. Berikut
bukti langsungnya: misalkan $O$ \omterm{prop:g10:coordgeom:midpoint}{titik tengah} $[AB]$ dan $r = \frac{AB}{2}$. Maka
$\vect{MA} = \vect{MO} + \vect{OA}$ dan
$\vect{MB} = \vect{MO} + \vect{OB} = \vect{MO} - \vect{OA}$, sehingga
\[
\vect{MA}\cdot\vect{MB}
= \norm{\vect{MO}}^2 - \norm{\vect{OA}}^2 = MO^2 - r^2,
\]
yang lenyap tepat ketika $MO = r$, yaitu ketika $M$ berada pada lingkaran
berpusat $O$ dan berjari-jari $r$.
\end{proof}

\begin{proposition}[Vektor normal sebuah garis]\label{prop:g11:scal:normal}
\omterm{def:g11:vect:vector}{Vektor} $\vec n\,(a, b)$ \omterm{def:g11:scal:orthogonal}{ortogonal} terhadap setiap \omterm{def:g11:vect:direction}{vektor arah}
garis $d : ax + by + c = 0$; kita katakan $\vec n$ adalah \emph{vektor
normal}\index{vektor normal} $d$. Dua garis saling tegak lurus tepat
ketika \omterm{def:g11:vect:vector}{vektor} normalnya (atau \omterm{def:g11:vect:direction}{vektor arahnya}) \omterm{def:g11:scal:orthogonal}{ortogonal}.
\end{proposition}

\begin{proof}
Sebuah \omterm{def:g11:vect:direction}{vektor arah} $d$ adalah $\vec u\,(-b, a)$
(\cref{thm:g11:vect:cartesian}), dan
$\vec n \cdot \vec u = a(-b) + b(a) = 0$.
\end{proof}

\begin{method}[Memakai vektor normal]\label{met:g11:scal:normal}
\begin{itemize}
  \item Garis melalui $A$ dengan normal $\vec n\,(a,b)$: tulislah
        $\vect{AM} \cdot \vec n = 0$, yang menjabar menjadi $ax + by + c = 0$
        dengan $c$ ditentukan oleh $A$.
  \item Ketegaklurusan $d: ax+by+c = 0$ dan $d': a'x+b'y+c' = 0$:
        periksalah $aa' + bb' = 0$.
\end{itemize}
\end{method}

\begin{example}\label{ex:g11:scal:perpline}
Garis melalui $P(2, -1)$ yang tegak lurus $d : 3x - y + 4 = 0$: sebuah
\omterm{def:g11:vect:direction}{vektor arah} $d$ adalah $(1, 3)$, yang berperan sebagai \omterm{def:g11:vect:vector}{vektor} \emph{normal}
garis tegak lurusnya. \omterm{def:g10:algebra:equation}{Persamaannya}: $x + 3y + c = 0$ dengan
$2 - 3 + c = 0$, jadi $c = 1$: garisnya adalah $x + 3y + 1 = 0$.
\end{example}

\begin{example}[Persamaan sebuah lingkaran]\label{ex:g11:scal:circle}
Lingkaran berpusat $\Omega(1, 2)$ dan berjari-jari $3$ adalah himpunan titik
$M(x,y)$ dengan $\Omega M^2 = 9$:
\[
(x - 1)^2 + (y - 2)^2 = 9 .
\]
Sebaliknya, $x^2 + y^2 - 2x - 4y - 4 = 0$ ditulis ulang, dengan melengkapkan
kuadratnya (\cref{ch:g11:quad}), menjadi $(x-1)^2 + (y-2)^2 = 9$: yaitu
lingkaran yang sama.
\end{example}

\section{Latihan}

\begin{exercise}[$\star$]\label{exo:g11:scal:1}
Hitunglah $\vec u \cdot \vec v$ untuk
\[
\vec u\,(2, 5),\ \vec v\,(3, -1); \qquad
\vec u\,(1, -4),\ \vec v\,(8, 2); \qquad
\norm{\vec u} = 3,\ \norm{\vec v} = 4,\ \theta = \frac{\pi}{3} .
\]
\end{exercise}

\begin{exercise}[$\star$]\label{exo:g11:scal:2}
$ABC$ adalah segitiga sama sisi bersisi $6$. Hitunglah
$\vect{AB} \cdot \vect{AC}$ dan $\vect{AB} \cdot \vect{BC}$.
\end{exercise}

\begin{exercise}[$\star$]\label{exo:g11:scal:3}
Untuk nilai $t$ yang mana $\vec u\,(3, t)$ dan $\vec v\,(t - 8, 1)$
\omterm{def:g11:scal:orthogonal}{ortogonal}? Lalu untuk nilai $t$ yang mana $\vec u\,(t, 2)$ \omterm{def:g11:scal:orthogonal}{ortogonal}
terhadap dirinya dikurangi $\vec v\,(4, 2)$, yaitu $\vec u \perp (\vec u - \vec v)$?
\end{exercise}

\begin{exercise}[$\star$]\label{exo:g11:scal:4}
Hitunglah sudut $\theta$ antara $\vec u\,(1, 2)$ dan $\vec v\,(3, 1)$
(berikan $\cos\theta$ secara eksak, lalu $\theta$ sampai derajat terdekat).
\end{exercise}

\begin{exercise}[$\star\star$]\label{exo:g11:scal:5}
Sebuah segitiga mempunyai sisi $b = 4$, $c = 6$ dan sudut
$\widehat A = \frac{2\pi}{3}$ di antaranya. Hitunglah sisi ketiganya, $a$.
Lalu, pada segitiga bersisi $a = 7$, $b = 5$, $c = 4$, hitunglah
$\cos\widehat A$ dan putuskan apakah $\widehat A$ lancip atau tumpul.
\end{exercise}

\begin{exercise}[$\star\star$]\label{exo:g11:scal:6}
Misalkan $A(1, 1)$, $B(5, -1)$ dan $C(3, 5)$. Hitunglah
$\vect{AB} \cdot \vect{AC}$; apakah sudut $\widehat A$ siku-siku? Hitunglah
ketiga panjang sisinya lalu periksa jawabanmu dengan aturan \omterm{def:g11:trigo:cossin}{kosinus}.
\end{exercise}

\begin{exercise}[$\star\star$]\label{exo:g11:scal:7}
Berikan \omterm{def:g10:algebra:equation}{persamaan} lingkaran berpusat $\Omega(-2, 3)$ dan berjari-jari $5$.
Lalu tentukan pusat dan jari-jari lingkaran
\[
x^2 + y^2 - 6x + 2y - 6 = 0 .
\]
\end{exercise}

\begin{exercise}[$\star\star$]\label{exo:g11:scal:8}
Misalkan $A(-3, 0)$ dan $B(0, 2)$. Dengan \cref{prop:g11:scal:diameter},
carilah \omterm{def:g10:algebra:equation}{persamaan} lingkaran berdiameter $[AB]$, lalu periksalah bahwa titik
asal $O(0,0)$ terletak padanya.
\end{exercise}

\begin{exercise}[$\star\star$]\label{exo:g11:scal:9}
Carilah \omterm{def:g10:algebra:equation}{persamaan} garis melalui $P(1, 3)$ yang tegak lurus terhadap
$d : 2x + y - 7 = 0$, dan \omterm{def:g10:coordgeom:system}{koordinat} kaki $H$ garis tegak
lurusnya (yaitu perpotongan kedua garis itu). Simpulkan jarak
dari $P$ ke $d$.
\end{exercise}

\begin{exercise}[$\star\star$]\label{exo:g11:scal:10}
Dengan memakai $\norm{\vec u + \vec v}^2 = \norm{\vec u}^2 + 2\vec u\cdot\vec v +
\norm{\vec v}^2$ dan padanannya untuk $\vec u - \vec v$, buktikan
\emph{identitas jajargenjang}:
\[
\norm{\vec u + \vec v}^2 + \norm{\vec u - \vec v}^2
= 2\norm{\vec u}^2 + 2\norm{\vec v}^2 ,
\]
lalu tafsirkan pada sebuah jajargenjang (diagonal terhadap sisinya).
\end{exercise}

\begin{exercise}[$\star\star\star$]\label{exo:g11:scal:11}
Misalkan $A$ dan $B$ dua titik dengan $AB = 4$, dan $I$ \omterm{prop:g10:coordgeom:midpoint}{titik tengah}
$[AB]$.
\begin{enumerate}
  \item Tunjukkan bahwa untuk setiap titik $M$ berlaku
        $\vect{MA}\cdot\vect{MB} = MI^2 - 4$.
        (Sisipkan $I$: $\vect{MA} = \vect{MI} + \vect{IA}$.)
  \item Simpulkan himpunan titik $M$ yang memenuhi
        $\vect{MA}\cdot\vect{MB} = 12$.
\end{enumerate}
\end{exercise}

\section{Soal: Dua hadiah hasil kali skalar}

\begin{problem}[{Soal akhir pekan --- rumus penjumlahan sudut jatuh
dari satu hasil kali titik, jarak ke sebuah garis dari yang lain,
dan ketiga garis tinggi bertemu sebagai bonus}]
\label{pb:g11:scal:1}
\omterm{def:g11:scal:dot}{Hasil kali skalar} tampak seperti pembukuan --- kalikan, jumlahkan ---
namun ia menyerahkan, hampir tanpa ongkos, dua harta yang
bertahan melawan berabad-abad perburuan sudut: \emph{rumus penjumlahan sudut}
dalam trigonometri (hitunglah $\cos 15^\circ$ secara eksak!) dan
\emph{jarak dari sebuah titik ke sebuah garis} dalam satu pecahan yang rapi.
Sebagai hadiah perpisahan, ia membuktikan bahwa ketiga garis tinggi setiap
segitiga melalui satu titik. Keempat wajah pada
\cref{thm:g11:scal:faces} akan bertugas semuanya.

\medskip
\noindent\textbf{Bagian I --- Kefasihan.}
\begin{enumerate}
  \item Untuk $\vec u(3, 4)$ dan $\vec v(-1, 2)$: hitunglah
        $\vec u \cdot \vec v$, kedua \omterm{def:g11:scal:dot}{normanya}, dan sudut
        di antara kedua \omterm{def:g11:vect:vector}{vektornya} (sampai sepersepuluh derajat).
  \item Carilah $k$ agar $(2, k)$ dan $(3, -4)$ \omterm{def:g11:scal:orthogonal}{ortogonal}
        (\cref{def:g11:scal:orthogonal}).
  \item Wajah fisikanya: sebuah kereta luncur ditarik $10$ m oleh gaya $50$ N
        pada $60^\circ$ terhadap tanah. Hitunglah usahanya,
        $W = \vec F \cdot \vec d$.
  \item Aturan \omterm{def:g11:trigo:cossin}{kosinus} (\cref{thm:g11:scal:cosines}): dua sisi
        sepanjang $5$ dan $7$ mengapit sudut $60^\circ$; carilah
        sisi ketiganya secara eksak.
  \item Kebalikannya: sebuah segitiga bersisi $4$, $6$, $8$. Hitunglah
        \omterm{def:g11:trigo:cossin}{kosinus} sudut yang berhadapan dengan sisi $8$, golongkan
        sudutnya, lalu berikan besarnya sampai sepersepuluh derajat.
\end{enumerate}

\medskip
\noindent\textbf{Bagian II --- Hadiah pertama: rumus penjumlahan
sudut.}
Misalkan $\vec u = (\cos a, \sin a)$ dan $\vec v = (\cos b, \sin b)$:
dua \omterm{def:g11:vect:vector}{vektor} satuan dengan sudut $a$ dan $b$.
\begin{enumerate}[resume]
  \item Hitunglah $\vec u \cdot \vec v$ dua kali --- sekali dengan
        \omterm{def:g10:coordgeom:system}{koordinat}, sekali dengan \omterm{def:g11:scal:dot}{norma} dan sudut di antara kedua
        \omterm{def:g11:vect:vector}{vektornya} --- lalu simpulkan rumus induknya:
        \[
        \cos(a - b) = \cos a \cos b + \sin a \sin b .
        \]
  \item Substitusikan $-b$ menggantikan $b$ (ingat paritas \omterm{def:g11:trigo:cossin}{sinus} dan \omterm{def:g11:trigo:cossin}{kosinus}!)
        untuk memperoleh $\cos(a + b)$.
  \item Perolehlah $\sin(a + b)$ dari identitas pelengkapnya,
        $\sin x = \cos\left(\frac\pi2 - x\right)$, yang diterapkan pada
        $x = a + b$.
  \item Nilai eksak akhirnya: hitunglah $\cos 15^\circ$ (sebagai
        $\cos(45^\circ - 30^\circ)$) dan $\sin 75^\circ$. Periksalah
        keduanya secara numeris.
  \item Ambil $b = a$ untuk menurunkan rumus sudut ganda:
        $\cos 2a$ (dalam ketiga bentuknya) dan $\sin 2a$.
  \item Dari $\cos 2a = 1 - 2\sin^2 a$, hitunglah
        $\sin^2 15^\circ$, lalu $\sin 15^\circ$ sebagai \omterm{def:g11:quad:discriminant}{akar}
        bersarang --- lalu periksalah bahwa kuadratnya cocok dengan
        nilai $\frac{\sqrt6 - \sqrt2}{4}$ yang diberikan metode
        pertanyaan 9.
\end{enumerate}

\medskip
\noindent\textbf{Bagian III --- Hadiah kedua: jarak ke sebuah
garis.}
\begin{enumerate}[resume]
  \item Misalkan $d$ garis $ax + by + c = 0$, dengan \omterm{def:g11:vect:vector}{vektor} normal
        $\vec n(a, b)$ (\cref{prop:g11:scal:normal}),
        dan $P(x_0, y_0)$ sebuah titik. Dengan mengambil sebarang titik $A$
        pada $d$, jelaskan mengapa jarak dari $P$ ke $d$ adalah
        $\dfrac{\abs{\vect{AP} \cdot \vec n}}{\norm{\vec n}}$,
        lalu simpulkan rumus
        \[
        \operatorname{dist}(P, d)
        = \frac{\abs{a x_0 + b y_0 + c}}{\sqrt{a^2 + b^2}} .
        \]
        Hitunglah jarak dari $P(5, 6)$ ke
        $3x + 4y - 12 = 0$.
  \item Tulislah \omterm{def:g10:algebra:equation}{persamaan} lingkaran berpusat $(5, 6)$
        yang menyinggung garis itu.
  \item Periksa silang rumusnya: hitunglah jarak dari
        titik asal ke $3x + 4y - 12 = 0$, lalu benarkan dengan
        menghitung kaki garis tegak lurus dari titik asalnya
        secara eksplisit.
  \item Segitiga $A(0,0)$, $B(6,0)$, $C(2,5)$: hitunglah luasnya
        dengan alas $[AB]$; lalu hitunglah jarak dari
        $B$ ke garis $(AC)$ dan periksalah bahwa jarak itu menghasilkan luas
        yang sama.
  \item Lingkaran berpusat $(2, 1)$ dan berjari-jari $2$, terhadap
        garis $x + y - 6 = 0$: hitunglah jarak pusatnya
        ke garis itu lalu golongkan kedudukannya (memotong,
        menyinggung, atau terpisah).
\end{enumerate}

\medskip
\noindent\textbf{Bagian IV --- Bonusnya: garis tinggi bertemu di satu titik.}
\begin{enumerate}[resume]
  \item Perumumlah \cref{exo:g11:scal:11}: untuk sebarang dua titik
        $A$, $B$ dengan \omterm{prop:g10:coordgeom:midpoint}{titik tengah} $I$, buktikan
        \[
        \vect{MA} \cdot \vect{MB} = MI^2 - \frac{AB^2}{4}
        \quad\text{untuk setiap titik } M .
        \]
  \item Simpulkan \cref{prop:g11:scal:diameter} dalam satu baris:
        himpunan titik $M$ dengan
        $\vect{MA} \cdot \vect{MB} = 0$ adalah lingkaran
        berdiameter $[AB]$.
  \item Garis tingginya: misalkan $H$ perpotongan
        garis tinggi dari $A$ dan dari $B$ pada segitiga $ABC$
        (sehingga $\vect{HA} \cdot \vect{BC} = 0$ dan
        $\vect{HB} \cdot \vect{CA} = 0$). Buktikan identitas
        \[
        \vect{HA} \cdot \vect{BC} + \vect{HB} \cdot \vect{CA}
        + \vect{HC} \cdot \vect{AB} = 0
        \]
        --- jabarkan semuanya dari $H$ dengan Chasles --- lalu
        simpulkan bahwa $H$ juga terletak pada garis tinggi dari $C$:
        jadi ketiga garis tingginya berpotongan di satu titik (yaitu
        \emph{titik tinggi}\index{titik tinggi}).
  \item Penutup: empat wajah, tiga piala. Wajah yang mana pada
        \cref{thm:g11:scal:faces} yang mengusung rumus penjumlahan sudut,
        yang mana rumus jaraknya, yang mana \omterm{pb:g11:scal:1}{titik tingginya} ---
        dan mengapa ``hitunglah besaran yang sama dua kali'' menjadi semboyan
        seluruh bab ini?

\end{enumerate}
\end{problem}
